\documentclass[a4paper,12pt]{article}
\usepackage[T2A]{fontenc}
\usepackage[utf8]{inputenc}
\usepackage[russian]{babel}
\usepackage[margin=1in]{geometry}
\usepackage{hyperref}
\usepackage{listings}
\usepackage{xcolor}
\usepackage{titlesec}
\usepackage{fancyvrb}

\definecolor{codebg}{rgb}{0.95,0.95,0.95}
\lstset{
	backgroundcolor=\color{codebg},
	basicstyle=\ttfamily\small,
	breaklines=true,
	language=Ruby,
	showstringspaces=false,
	tabsize=2
}

\title{OpenAlexWrap\\\large{Ruby-обёртка для OpenAlex API}}
\author{}
\date{}

\begin{document}
	
	\maketitle
	
	\section{Введение}
	
	OpenAlexWrap — это Ruby-обёртка для \href{https://openalex.org/}{OpenAlex API}, предоставляющая чистый и интуитивно понятный интерфейс для доступа к метаданным научных публикаций, включая работы, авторов, источники и темы.
	
	\section{Возможности}
	
	\begin{itemize}
		\item \textbf{Поиск работ, авторов, источников и тем} с гибкой фильтрацией
		\item \textbf{Получение наиболее цитируемых/продуктивных сущностей} с настраиваемыми лимитами
		\item \textbf{Получение сущностей по идентификаторам} (DOI, ORCID, ISSN и т.д.)
		\item \textbf{Простота использования} с Ruby-подобным синтаксисом
	\end{itemize}
	
	\section{Установка}
	
	Добавьте эту строку в Gemfile вашего приложения:
	
	\begin{lstlisting}
		gem 'open_alex_wrap'
	\end{lstlisting}
	
	Затем выполните:
	
	\begin{lstlisting}[language=bash]
		$ bundle install
	\end{lstlisting}
	
	Или установите самостоятельно:
	
	\begin{lstlisting}[language=bash]
		$ gem install open_alex_wrap
	\end{lstlisting}
	
	\section{Использование}
	
	\subsection{Инициализация клиента}
	
	\begin{lstlisting}
		require 'open_alex_wrap'
		
		client = OpenAlexWrap::Client.new(
		email: 'your-email@example.com',  # Required
		api_key: 'your-api-key',           # Optional, but recommended                                 for higher rate limits
		timeout: 10                         # Optional, default is 10                                     seconds
		)
	\end{lstlisting}
	
	\subsection{Работа с работами (Works)}
	
	\begin{lstlisting}
		# Get most cited works of all time
		most_cited = client.most_cited_works(limit: 10)
		
		# Get most cited works from 2023
		most_cited_2023 = client.most_cited_works(year: 2023, limit: 20)
		
		# Search for works by keywords
		ai_works = client.search_works('artificial intelligence', limit: 50)
		
		# Search with filters
		open_access_ai = client.search_works('machine learning', open_access: true, limit: 30)
		
		# Get a work by DOI
		work = client.work_by_doi('10.1038/s41586-020-2649-2')
		
		# Get works by author
		author_works = client.works_by_author('A123456789', limit: 100)
		
		# Get works from a specific source (journal)
		nature_works = client.works_by_source('S123456789', year: 2024, limit: 50)
		
		# Get open access works
		oa_works = client.open_access_works(year: 2023, limit: 100)
		
		# Get highly cited works
		highly_cited = client.highly_cited_works(min_citations: 500, year: 2022)
		
		# Iterate through all works matching criteria (cursor-based pagination)
		count = 0
		client.each_work(filter: { publication_year: 2023 }) do |work|
		count += 1
		puts work['title']
		end
		puts "Total works: #{count}"
	\end{lstlisting}
	
	\subsection{Работа с авторами (Authors)}
	
	\begin{lstlisting}
		# Get most cited authors
		top_authors = client.most_cited_authors(limit: 50)
		
		# Get most prolific authors (by number of works)
		prolific = client.most_prolific_authors(limit: 30)
		
		# Search authors by name
		smith_authors = client.search_authors('John Smith', limit: 20)
		
		# Get authors from a specific institution
		harvard_authors = client.authors_by_institution('I123456789', limit: 100)
		
		# Get author by ORCID
		author = client.author_by_orcid('0000-0000-0000-0000')
		
		# Get authors by h-index (client-side sorting)
		top_hindex = client.top_authors_by_hindex(limit: 100)
	\end{lstlisting}
	
	\subsection{Работа с источниками (Sources)}
	
	\begin{lstlisting}
		# Get most cited sources
		top_journals = client.most_cited_sources(limit: 50)
		
		# Search sources by name
		science_journals = client.search_sources('science', limit: 30)
		
		# Get open access sources (from DOAJ)
		oa_journals = client.open_access_sources(limit: 50)
		
		# Get sources by publisher
		elsevier_sources = client.sources_by_publisher_id('P4310320068', limit: 30)
		
		# Get sources by country
		us_journals = client.sources_by_country('US', limit: 100)
		
		# Get source by ISSN
		nature = client.source_by_issn('1476-4687')
		
		# Get most prolific sources
		most_papers = client.most_prolific_sources(limit: 50)
	\end{lstlisting}
	
	\subsection{Работа с темами (Topics)}
	
	\begin{lstlisting}
		# Get a topic by ID
		topic = client.topic('T11636')
		
		# Get most popular topics (by number of works)
		popular = client.most_popular_topics(limit: 50)
		
		# Search topics by name
		machine_learning = client.search_topics('machine learning', limit: 30)
		
		# Get works by topic
		topic_works = client.works_by_topic('T11636', limit: 100)
		
		# Get topics by domain (Physical Sciences, Life Sciences, etc.)
		physical_sciences = client.topics_by_domain('D1', limit: 50)
		
		# Get topic hierarchy (domain  field  subfield  topic)
		hierarchy = client.topic_hierarchy('T11636')
		hierarchy.each do |level|
		puts "#{level[:level]}: #{level[:name]}"
		end
		
		# Get topic statistics
		stats = client.topic_stats('T11636')
		puts "Works: #{stats[:works_count]}"
		puts "Citations: #{stats[:cited_by_count]}"
		puts "Domain: #{stats[:domain]}"
		
		# Iterate through all topics matching criteria
		client.each_topic(filter: { works_count: '>1000' }) do |topic|
		puts "#{topic['display_name']}: #{topic['works_count']} works"
		end
	\end{lstlisting}
	
	\subsection{Расширенная фильтрация}
	
	Метод \texttt{query} поддерживает различные форматы фильтров:
	
	\begin{lstlisting}
		# Hash format (recommended)
		client.query(:work, filter: { publication_year: 2023, open_access: true })
		
		# String format
		client.query(:work, filter: 'publication_year:2023,open_access:true')
	\end{lstlisting}
	
	\subsection{Сортировка}
	
	\begin{lstlisting}
		# Hash format
		client.query(:work, sort: { cited_by_count: :desc, publication_date: :asc })
		
		# String format
		client.query(:work, sort: 'cited_by_count:desc,publication_date:asc')
	\end{lstlisting}
	
	\subsection{Выбор конкретных полей}
	
	Повысьте производительность, выбирая только необходимые поля:
	
	\begin{lstlisting}
		# Array format
		client.query(:work, select: ['title', 'doi', 'publication_year'])
		
		# String format
		client.query(:work, select: 'title,doi,publication_year')
	\end{lstlisting}
	
	\subsection{Пагинация с курсорами}
	
	\begin{lstlisting}
		# Use the each_work method for automatic cursor pagination
		client.each_work(filter: { publication_year: 2023 }) do |work|
		# Process each work
		end
		
		# Or handle cursor manually
		response = client.query(:work, filter: { publication_year: 2023 }, cursor: '*')
		while response['results']&.any?
		response['results'].each { |work| puts work['title'] }
		response = client.query(:work, cursor: response['meta']['next_cursor'])
		break if response['meta']['next_cursor'].nil?
		end
	\end{lstlisting}
	
	\section{Доступные типы сущностей}
	
	\begin{itemize}
		\item \textbf{Works} (\texttt{:work}) — научные статьи, публикации, книги и т.д.
		\item \textbf{Authors} (\texttt{:author}) — исследователи и ученые
		\item \textbf{Sources} (\texttt{:source}) — журналы
		\item \textbf{Topics} (\texttt{:topic}) — исследовательские темы и концепции
	\end{itemize}
	
	
	\section{Лимитирование запросов}
	
	OpenAlex API имеет лимиты на количество запросов.
	
	
	\section{Ссылки}
	
	\begin{itemize}
		\item \href{https://docs.openalex.org/}{OpenAlex API Documentation}
		\item \href{https://github.com/ZubtsoVA/open_alex_wrap}{GitHub Repository}
	\end{itemize}
	
\end{document}